\section{Setting}

A feedforward MLP of depth $L$ and width $N$ maps $x\in\R^{D}$ to a scalar by
\begin{align}
  h^{(0)} &= x, \nonumber\\
  h^{(\ell+1)} &= N^{-a_{\ell+1}} W^{(\ell+1)}\,\ph\!\big(h^{(\ell)}\big),
  \qquad \ell=0,\ldots,L-2, \nonumber\\
  f &= \gamma^{-1} N^{-a_L}\, w\cdot\ph\!\big(h^{(L-1)}\big).
  \label{eq:mlp-schema}
\end{align}
The activation $\ph$ is applied elementwise. Weights are drawn $W^{(\ell)}_{ij}\sim\mathcal{N}(0,N^{-b_\ell})$, and the training set is $\data=\{(x_\mu,y_\mu)\}_{\mu=1}^{P}$. Two kernels appear throughout:
\begin{equation}
  \Phi^{(\ell)}_{\mu\nu}
  \equiv
  N^{-1}\,\ph\!\big(h^{(\ell)}_\mu\big)\cdot\ph\!\big(h^{(\ell)}_\nu\big),
  \qquad
  K_{\mu\nu}(t)
  \equiv
  \partial_\theta f(x_\mu;t)\cdot\partial_\theta f(x_\nu;t).
  \label{eq:two-kernels}
\end{equation}
The first is the \emph{feature kernel} of layer $\ell$. The second is the \emph{neural tangent kernel}. In the limits treated here both concentrate, and the $O(N)$ weights are replaced by $O(P^2)$ kernel entries.

The questions are: which functions~\eqref{eq:mlp-schema} can represent; how $f(\,\cdot\,;\theta(t))$ moves under gradient flow; and what the population risk is after training. The answers depend on how $(a_\ell,b_\ell,\gamma)$ scale with $N$, and on the ratio $L/N$.
